\chapter{RANCANGAN PENELITIAN}
\label{chap:rancangan}
% Rancangan kegiatan untuk menjawab pertanyaan penelitian dan mencapai tujuan
% penelitian, termasuk metode penelitian dan jadwal pelaksanaan yang rasional.

Bab ini menjabarkan rancangan kegiatan yang akan dilakukan untuk memperoleh jawaban terhadap pertanyaan penelitian dan mencapai tujuan penelitian, termasuk metode penelitian serta jadwal pelaksanaan penelitian.

\section{Metode Penelitian}
Jelaskan metode atau paradigma penelitian yang digunakan (misalnya: penelitian eksperimen, \textit{design science research}, studi kasus, kuantitatif, kualitatif, atau campuran), alasan pemilihannya, serta kaitannya dengan pertanyaan dan tujuan penelitian.

\section{Tahapan Penelitian}
\label{sec:tahapan}
Uraikan tahapan penelitian secara berurutan dan sistematis. Setiap tahapan harus memiliki masukan, proses, dan luaran yang jelas. Gambarkan alur tahapan dalam diagram seperti Gambar~\ref{fig:tahapan}, lalu jelaskan setiap tahapannya.

\begin{figure}[H]
  \centering
  \begin{tikzpicture}[node distance=2.2cm]
  \node (studi)    [process]                          {Studi Literatur};
  \node (data)     [process, below of=studi]          {Pengumpulan \& Persiapan Data};
  \node (rancang)  [process, below of=data]           {Perancangan Solusi};
  \node (impl)     [process, below of=rancang]        {Implementasi \& Eksperimen};
  \node (valid)    [decision, below of=impl, yshift=-0.5cm] {Target tercapai?};
  \node (revisi)   [process2, right of=valid, xshift=2.5cm] {Perbaikan};
  \node (evaluasi) [process, below of=valid, yshift=-0.5cm] {Evaluasi \& Penarikan Kesimpulan};

  \draw [arrow] (studi)   -- (data);
  \draw [arrow] (data)    -- (rancang);
  \draw [arrow] (rancang) -- (impl);
  \draw [arrow] (impl)    -- (valid);
  \draw [arrow] (valid)   -- node[anchor=south] {Tidak} (revisi);
  \draw [arrow] (revisi)  |- (impl);
  \draw [arrow] (valid)   -- node[right] {Ya} (evaluasi);
  \end{tikzpicture}
  \caption{Tahapan Penelitian}
  \label{fig:tahapan}
\end{figure}

\section{Pemetaan Pertanyaan Penelitian}
Tunjukkan bagaimana setiap pertanyaan penelitian akan dijawab: kegiatan yang dilakukan, luaran yang dihasilkan, dan cara mengevaluasinya, seperti pada Tabel~\ref{tab:pemetaan_rq}.

\begin{table}[htbp]
  \centering
  \caption{Pemetaan Pertanyaan Penelitian terhadap Kegiatan Penelitian}
  \label{tab:pemetaan_rq}
  \begin{tabularx}{\textwidth}{|>{\centering\arraybackslash}p{0.08\textwidth}|*{3}{>{\RaggedRight\arraybackslash}X|}}
  \hline
  \textbf{RQ} & \textbf{Kegiatan} & \textbf{Luaran} & \textbf{Cara Evaluasi} \\ \hline
  RQ1 & [Kegiatan untuk menjawab RQ1] & [Artefak / model / temuan] & [Metrik / teknik evaluasi] \\ \hline
  RQ2 & [Kegiatan untuk menjawab RQ2] & [Artefak / model / temuan] & [Metrik / teknik evaluasi] \\ \hline
  \end{tabularx}
\end{table}

\section{Rancangan Eksperimen dan Evaluasi}
Jelaskan skenario eksperimen yang akan dilakukan (misalnya: variasi parameter, pembanding/\textit{baseline}, pembagian data, jumlah pengulangan) serta prosedur evaluasi hasil menggunakan metrik yang telah diuraikan pada Bab~II.

\section{Alat dan Bahan Penelitian}
Cantumkan perangkat keras (\textit{hardware}), perangkat lunak (\textit{software}), bahasa pemrograman, \textit{framework}, atau dataset yang akan digunakan, beserta versi dan alasan pemilihannya.

\begin{table}[htbp]
  \begin{center}
    \caption{Alat dan Bahan Penelitian}
    \label{tab:tools}
    \begin{tabularx}{\textwidth}{|>{\centering\arraybackslash}p{0.06\textwidth}|>{\RaggedRight\arraybackslash}X|>{\RaggedRight\arraybackslash}X|>{\RaggedRight\arraybackslash}X|}
    \hline
    \textbf{No.} & \textbf{Alat / Bahan} & \textbf{Spesifikasi / Keterangan} & \textbf{Fungsi dalam Penelitian} \\ \hline
    1 & [Perangkat Lunak / Framework] & [Nama, Versi] & [Fungsi] \\ \hline
    2 & [Bahasa Pemrograman] & [Nama, Versi] & [Fungsi] \\ \hline
    3 & [Basis Data] & [Nama, Versi] & [Fungsi] \\ \hline
    4 & [Dataset / Sumber Data] & [Deskripsi singkat] & [Fungsi] \\ \hline
    \end{tabularx}
  \end{center}
\end{table}

\section{Jadwal Pelaksanaan Penelitian}
Jadwal pelaksanaan penelitian disajikan pada Tabel~\ref{tab:jadwal}. Susun jadwal secara rasional dan sesuaikan aktivitas serta periode waktunya dengan tahapan penelitian pada Subbab~\ref{sec:tahapan}. Jadwal ini mencakup keseluruhan siklus penelitian, mulai dari perumusan masalah hingga Sidang Tugas Akhir.

\begin{table}[htbp]
  \caption{Jadwal Pelaksanaan Penelitian}
  \label{tab:jadwal}
  \centering
  \adjustbox{max width=\textwidth}{%
    \begin{tabular}{|c|>{\RaggedRight\arraybackslash}p{0.22\linewidth}|c|c|c|c|c|c|c|c|c|c|c|}
    \hline
    \textbf{No.} & \textbf{Kegiatan}
      & \textbf{\begin{tabular}[c]{@{}c@{}}Agt\\2026\end{tabular}}
      & \textbf{\begin{tabular}[c]{@{}c@{}}Sep\\2026\end{tabular}}
      & \textbf{\begin{tabular}[c]{@{}c@{}}Okt\\2026\end{tabular}}
      & \textbf{\begin{tabular}[c]{@{}c@{}}Nov\\2026\end{tabular}}
      & \textbf{\begin{tabular}[c]{@{}c@{}}Des\\2026\end{tabular}}
      & \textbf{\begin{tabular}[c]{@{}c@{}}Jan\\2027\end{tabular}}
      & \textbf{\begin{tabular}[c]{@{}c@{}}Feb\\2027\end{tabular}}
      & \textbf{\begin{tabular}[c]{@{}c@{}}Mar\\2027\end{tabular}}
      & \textbf{\begin{tabular}[c]{@{}c@{}}Apr\\2027\end{tabular}}
      & \textbf{\begin{tabular}[c]{@{}c@{}}Mei\\2027\end{tabular}}
      & \textbf{\begin{tabular}[c]{@{}c@{}}Jun\\2027\end{tabular}} \\
    \hline
    1 & Identifikasi \& Perumusan Masalah
      & \cellcolor{green!25} & \cellcolor{green!25} & & & & & & & & & \\
    \hline
    2 & Studi Literatur
      & & \cellcolor{green!25} & \cellcolor{green!25} & \cellcolor{green!25} & & & & & & & \\
    \hline
    3 & Pengumpulan Data
      & & & \cellcolor{green!25} & \cellcolor{green!25} & & & & & & & \\
    \hline
    4 & Analisis Kebutuhan
      & & & \cellcolor{green!25} & \cellcolor{green!25} & & & & & & & \\
    \hline
    5 & Perancangan (\textit{Design})
      & & & & \cellcolor{green!25} & \cellcolor{green!25} & & & & & & \\
    \hline
    6 & \textbf{Seminar Proposal}
      & & & & \cellcolor{orange!50} & & & & & & & \\
    \hline
    7 & Implementasi
      & & & & & \cellcolor{green!25} & \cellcolor{green!25} & \cellcolor{green!25} & & & & \\
    \hline
    8 & \textbf{Seminar TA-1}
      & & & & & \cellcolor{orange!50} & & & & & & \\
    \hline
    9 & Pengujian \& Evaluasi
      & & & & & & \cellcolor{green!25} & \cellcolor{green!25} & \cellcolor{green!25} & & & \\
    \hline
    10 & Penarikan Kesimpulan \& Saran
      & & & & & & & & \cellcolor{green!25} & \cellcolor{green!25} & & \\
    \hline
    11 & \textbf{Seminar Hasil}
      & & & & & & & & & & \cellcolor{orange!50} & \\
    \hline
    12 & \textbf{Sidang Tugas Akhir}
      & & & & & & & & & & & \cellcolor{orange!50} \\
    \hline
    \end{tabular}%
  }
  \begin{flushleft}
    \small \textit{Keterangan:} \colorbox{green!25}{\phantom{X}} Kegiatan penelitian \quad \colorbox{orange!50}{\phantom{X}} Milestone / Seminar
  \end{flushleft}
\end{table}
